\section{Introducción}

Las pequeñas y medianas empresas (PyMEs) representan un pilar fundamental de las
economías latinoamericanas; sin embargo, su acceso a herramientas de inteligencia
de negocios y analítica de datos sigue siendo limitado. Según
\citet{tovar2011bi}, los principales obstáculos para la adopción de soluciones
de \textit{Business Intelligence} (BI) en este segmento empresarial son el costo
elevado de las plataformas, la falta de personal cualificado y la complejidad de
las integraciones con sistemas existentes.

Esta brecha tecnológica se profundiza cuando se considera que, a diferencia de
las grandes corporaciones, las PyMEs carecen de departamentos de TI dedicados
capaces de gestionar infraestructuras de datos complejas
\citep{cordero2020sistemas}. A esto se suma la incompatibilidad entre sistemas
heredados y las nuevas tendencias tecnológicas, lo cual constituye una barrera
estructural que dificulta la modernización del sector empresarial y limita sus
posibilidades de crecimiento. Asimismo, la complejidad percibida al intentar adoptar
nuevas tecnologías de la información termina siendo un obstáculo que desanima a
muchas empresas \citep{nurqamarani2021complejidad}. Las interfaces técnicas de
las herramientas de analítica requieren conocimientos especializados en bases de
datos, estadística y ciencia de datos que frecuentemente no están al alcance del
perfil medio de los trabajadores de una PyME, traduciéndose en una subutilización
de la información organizacional.

Incluso cuando las organizaciones logran implementar soluciones analíticas, la
falta de una cultura basada en datos y la escasez de competencias internas
generan un retorno de inversión insatisfactorio \citep{wamba2017desafios}. En
consecuencia, existe una necesidad urgente de soluciones de bajo código o
lenguaje natural que permitan democratizar el acceso a la información
estructurada.

Frente a este panorama, la Inteligencia Artificial Generativa emerge como una
tecnología catalizadora para resolver esta problemática. La capacidad de los
grandes modelos de lenguaje (LLMs) para procesar información y generar código
facilita la creación de sistemas analíticos más accesibles
\citep{barbon2024generative}. Proyecciones recientes indican que la
automatización mediante agentes de IA y la inteligencia de decisión son
tendencias clave para reducir la brecha analítica en el sector corporativo
\citep{gartner2025predicciones}.

Este trabajo presenta BIGENIA, una arquitectura de Inteligencia de Negocios
diseñada para mitigar estos obstáculos mediante el uso de lenguaje natural. El
sistema implementa un agente de \textit{Generative Business Intelligence}
(GenBI) capaz de interpretar consultas expresadas en lenguaje cotidiano,
convertirlas en código estructurado y desplegar los resultados mediante gráficos
interactivos de Plotly. Para validar el funcionamiento de esta arquitectura y su
precisión en la extracción de datos, se utiliza el conjunto de datos relacional
Chinook como entorno de prueba. De este modo, se busca proporcionar una interfaz
que elimine la necesidad de intermediarios técnicos o conocimientos avanzados en
SQL.

El presente artículo se estructura de la siguiente manera: la Sección 2 presenta
la justificación del trabajo. La Sección 3 detalla el marco teórico fundamental.
La Sección 4 revisa el estado del arte en IA aplicada a bases de datos.
Posteriormente, la Sección 5 describe de forma exhaustiva la arquitectura de la
aplicación, desglosando los agentes y componentes integrados. Finalmente, se
exponen las conclusiones obtenidas y se proponen futuras líneas de investigación.